\section{Evidence Matrix: cita y evidencia visual}

\subsection{Quote y mapeo temporal}
\label{sec:evidence-quotes}
Para garantizar que cada párrafo de la nota se alinee con el video original, presentamos los quotes clave en una tabla que incluye las palabras textuales del profesor, el timecode correspondiente y el lugar donde aparecen en la nota. Esto facilita la verificación durante el audit y también ayuda al future reader a converger rápidamente en ese argumento.

\begin{table}[H]
\footnotesize
\centering
\begin{tabular}{p{3cm}p{4cm}p{5cm}}
\toprule
Sección & Cita textual/hallazgo & Tiempo correspondiente \\
\midrule
Data & ``Dive into Kimi K2.5, este artículo tan extenso'', lo que indica que se aborda desde tres líneas: datos, algoritmo y arquitectura. & 00:00:20--00:00:40 \\
Algorithm & ``deg and recover es el ciclo natural de resincronización de GRM + PARL'', que nos recuerda no alarmarnos por retrocesos de corto plazo. & 00:18:40--00:19:05 \\
Architecture & ``Activate the Vision-Agentic Capability'' hace énfasis en considerar de manera conjunta la visión, el lenguaje y las herramientas. & 00:32:10--00:32:48 \\
Agent Swarm & ``S mean + max'' resume la forma de agregación de reward para múltiples agents en paralelo. & 00:35:20--00:36:50 \\
Infra & ``Critical steps son tanto guard band como compute barrier'', lo que explica la lógica subyacente del loop de evaluación. & 00:41:10--00:42:30 \\
\bottomrule
\end{tabular}
\caption{Sección y timecode correspondientes a los quotes clave, que ayudan a ubicar rápidamente el fragmento de video y el párrafo de la nota durante la revisión.}
\end{table}

\begin{knowledgebox}{La función de Evidence Matrix}
\begin{itemize}
    \item Vincula la cita textual, la sección y el tiempo, de manera que el future reviewer pueda verificar rápidamente el origen del contenido durante el audit.
    \item Antes de que sea necesario agregar una figure o un summary, ya deja claros el orden temporal y el tema, lo que evita que el contenido se desvíe.
    \item Es adecuado para hacer un check cruzado entre transcripts (`lecture09.srt`) y las figure footnote, lo que mejora la reproducibility.
\end{itemize}
\end{knowledgebox}

\subsection{Checklist de fotogramas/Slide}
Por el momento solo está disponible la imagen de la portada, por lo que preparamos además un checklist de fotogramas/slide que indica la posición que debe ocupar cada imagen, su timecode y su propósito (por ejemplo, highlight architecture, agent swarm, monitoring).

\begin{table}[H]
\footnotesize
\centering
\begin{tabular}{p{3cm}p{4cm}p{5cm}}
\toprule
Imagen objetivo & Intervalo de tiempo & Uso \\
\midrule
Portada / apertura del profesor & 00:00:05--00:00:13 & Explica las tres líneas de K2.5; se coloca en la introducción como anchor visual. \\
Agent loop demo & 00:32:10--00:32:48 & Muestra la figura compuesta de Visual Agency y agent swarm; se coloca en la sección de Architecture. \\
Critical steps dashboard & 00:41:10--00:42:30 & Se usa para la descripción de bandmark de Infrastructure, acompañado de una explicación en `figure`. \\
\bottomrule
\end{tabular}
\caption{Los available frames actuales y su intended placement; si más adelante hay slides, también se pueden agregar filas en la same table.}
\end{table}

\begin{importantbox}{Requisitos de registro de Frame/Slide}
Registra el local filename, el timecode, la sección a la que pertenece y el \verb|\protect\footnotemark| de cada figure, y organiza de una vez, dentro del flujo práctico, el punto técnico, el quote y el summary de esa imagen, para acelerar la revisión del audit.
\end{importantbox}

\subsection{Resumen de la sección}
Evidence Matrix nos permite alinear las tres dimensiones «quote/timecode/figure», de manera que cada hallazgo sea trazable; cuando los slides estén disponibles, también se pueden agregar en el mismo matrix, para mantener la consistencia del documento práctico.
